\documentclass[a4paper,12pt]{article}

\usepackage[T2A]{fontenc}
\usepackage[utf8]{inputenc}
\usepackage[russian]{babel}
\usepackage{amsmath,amssymb,amsthm,mathtools}
\usepackage{helvet}
\renewcommand{\familydefault}{\sfdefault}

% Стабильный sans-serif fallback для кириллицы при pdfLaTeX.
\DeclareFontFamily{T2A}{phv}{}
\DeclareFontShape{T2A}{phv}{m}{n}{<->ssub*cmss/m/n}{}
\DeclareFontShape{T2A}{phv}{b}{n}{<->ssub*cmss/bx/n}{}
\DeclareFontShape{T2A}{phv}{bx}{n}{<->ssub*cmss/bx/n}{}

\usepackage{icomma}
\usepackage[a4paper,left=20mm,right=20mm,top=20mm,bottom=22mm]{geometry}
\usepackage[nopatch=footnote]{microtype}
\usepackage{tikz}
\usepackage{tkz-euclide}
\usepackage{tabularx,booktabs}
\usepackage{enumitem}
\usepackage{hyperref}
\usepackage{parskip}
\usepackage{fancyhdr}
\usepackage{setspace}
\usepackage{xcolor}

\definecolor{accent}{HTML}{008080}
\definecolor{darkaccent}{HTML}{004D4D}
\definecolor{textgray}{HTML}{333333}
\definecolor{softgray}{HTML}{6B7280}
\definecolor{linegray}{HTML}{D7DEDE}

\hypersetup{
  colorlinks=true,
  linkcolor=darkaccent,
  urlcolor=darkaccent,
  pdfauthor={Лёвин Артём Александрович},
  pdftitle={Чек-лист: буквенные выражения, проценты, дроби и неравенства}
}

\onehalfspacing
\setlength{\parindent}{0pt}
\setlength{\headheight}{25pt}
\raggedbottom

\pagestyle{fancy}
\fancyhf{}
\fancyhead[L]{\small\textcolor{softgray}{Математика · 7 класс}}
\fancyhead[R]{\small\textcolor{softgray}{Кирилл Зиновьев}}
\fancyfoot[C]{\small\textcolor{softgray}{\thepage}}
\renewcommand{\headrulewidth}{0.3pt}
\renewcommand{\headrule}{%
  \hbox to\headwidth{%
    \color{linegray}\leaders\hrule height \headrulewidth\hfill
  }%
}

\newcommand{\sectiontitle}[1]{%
  \vspace{0.6em}%
  {\Large\bfseries\textcolor{darkaccent}{#1}}\par
  \vspace{0.25em}%
  {\color{linegray}\rule{\linewidth}{0.6pt}}\par
  \vspace{0.5em}%
}

\newcommand{\keyidea}[1]{%
  \textbf{\textcolor{darkaccent}{Ключевая идея.}} #1
}

\newcommand{\warning}[1]{%
  \textbf{\textcolor{accent}{Важно.}} #1
}

\setlist[enumerate,1]{
  label=\arabic*.,
  leftmargin=*,
  itemsep=0.65em,
  topsep=0.4em
}

\setlist[enumerate,2]{
  label=\arabic*),
  leftmargin=1.6em,
  itemsep=0.3em,
  topsep=0.2em
}

\begin{document}

% =========================================================
% ТИТУЛЬНЫЙ ЛИСТ
% =========================================================

\begin{titlepage}
\thispagestyle{empty}
\centering

\vspace*{2.5cm}

{\Large\textcolor{accent}{\bfseries ЧЕК-ЛИСТ}}\par

\vspace{0.8cm}

{\Huge\bfseries\textcolor{textgray}
{Буквенные выражения, проценты, дроби и простейшие неравенства}


\vspace{1.1cm}

{\large\textcolor{softgray}{
Подстановка · математические модели · общий знаменатель · числовая прямая
}}\par

\vfill

{\large\bfseries Лёвин Артём Александрович}\par

\vspace{0.2cm}

{\large эксклюзивно для Кирилла Зиновьева}\par

\vspace{0.6cm}

{\large \today}\par

\vfill

\begin{tikzpicture}[x=1cm,y=1cm]
  \draw[accent,line width=1pt]
    (-7,0)
    .. controls (-4.5,0.65) and (-2.2,-0.45) ..
    (0,0)
    .. controls (2.2,0.45) and (4.5,-0.65) ..
    (7,0);
\end{tikzpicture}

\vspace*{0.5cm}

\end{titlepage}

% =========================================================
% ОСНОВНОЙ ЧЕК-ЛИСТ
% =========================================================

\sectiontitle{1. Подстановка чисел в буквенное выражение}

Если значения букв известны, каждую букву заменяют соответствующим числом,
после чего выполняют вычисления по обычному порядку действий.

\[
23a-7b-2c,
\qquad
a=0{,}2,\quad b=5{,}4,\quad c=6{,}2.
\]

После подстановки:

\[
23\cdot 0{,}2
-
7\cdot 5{,}4
-
2\cdot 6{,}2.
\]

\keyidea{
Сначала аккуратно выполните подстановку, затем переходите к вычислениям.
Знаки чисел и знаки действий должны сохраниться.
}

\warning{
Если вместо буквы подставляется отрицательное число, заключайте его в скобки.
Например, при $x=-3$ выражение $5x^2$ превращается в
$5\cdot(-3)^2$.
}

% ---------------------------------------------------------

\sectiontitle{2. Если известно значение целого выражения}

Иногда значения отдельных переменных неизвестны, зато известно значение
некоторого выражения целиком. Такое выражение удобно воспринимать как один
готовый математический блок.

Если

\[
3xy-2z=4,
\]

то

\[
\frac{1}{3xy-2z}
=
\frac14
=
0{,}25.
\]

Значения $x$, $y$ и $z$ отдельно искать не требуется.

\keyidea{
Если в новом выражении полностью повторяется уже известная комбинация,
подставляйте её значение целиком.
}

\warning{
Для дроби знаменатель обязан быть ненулевым.
В данном примере условие выполнено, поскольку
$3xy-2z=4\ne0$.
}

% ---------------------------------------------------------

\sectiontitle{3. Перевод текста и рисунка в математическую модель}

Один из важнейших навыков в алгебре --- перевод условия задачи на язык формул.

Пусть дан квадрат со стороной $a$.
От него отрезали прямоугольную полосу со сторонами $a$ и $b$, где

\[
0<b<a.
\]

Оставшаяся фигура представляет собой прямоугольник со сторонами
$a$ и $a-b$.

\begin{center}
\begin{tikzpicture}[scale=0.9]

  \tkzDefPoint(0,0){A}
  \tkzDefPoint(5,0){B}
  \tkzDefPoint(5,5){C}
  \tkzDefPoint(0,5){D}

  \tkzDefPoint(0,1.7){E}
  \tkzDefPoint(5,1.7){F}

  \fill[accent!8] (0,1.7) rectangle (5,5);
  \fill[softgray!10] (0,0) rectangle (5,1.7);

  \tkzDrawPolygon[line width=1pt,color=accent](A,B,C,D)
  \tkzDrawSegment[line width=1pt,color=darkaccent](E,F)

  \tkzLabelPoints[below](A,B)
  \tkzLabelPoints[above](C,D)

  \draw (0,-0.45) -- (5,-0.45);

  \node[below,text=textgray] at (2.5,-0.45) {$a$};
  \node[left,text=textgray] at (0,2.5) {$a$};
  \node[left,text=textgray] at (0,0.85) {$b$};
  \node[left,text=darkaccent] at (0,3.35) {$a-b$};

  \node[text=darkaccent] at (2.5,3.35)
    {оставшаяся часть};

  \node[text=softgray] at (2.5,0.85)
    {отрезанная часть};

\end{tikzpicture}
\end{center}

Следовательно,

\[
S=a(a-b)=a^2-ab.
\]

\keyidea{
Сначала изобразите ситуацию, затем выразите неизвестный размер
через уже известные величины.
}

% ---------------------------------------------------------

\sectiontitle{4. Проценты: нахождение числа по его проценту}

Если $p\%$ некоторого числа $x$ равны числу $A$, то можно составить пропорцию

\[
\frac{p}{100}
=
\frac{A}{x}.
\]

Отсюда

\[
x=\frac{100A}{p},
\qquad
p\ne0.
\]

Например, $7\%$ некоторого числа равны $28$.

\[
\frac{7}{100}
=
\frac{28}{x}.
\]

Перемножаем крест-накрест:

\[
7x=28\cdot100.
\]

Следовательно,

\[
x=400.
\]

\keyidea{
За $100\%$ принимают всё число целиком.
Если целое неизвестно, обозначьте его через $x$.
}

% ---------------------------------------------------------

\sectiontitle{5. Общий знаменатель и сравнение отрицательных дробей}

Для дробей со знаменателями $4$, $6$ и $7$ удобно использовать

\[
\text{НОК}(4,6,7)=84.
\]

Например,

\[
-\frac54
=
-\frac{105}{84},
\qquad
-\frac76
=
-\frac{98}{84},
\qquad
-\frac87
=
-\frac{96}{84}.
\]

Поэтому

\[
-\frac54
<
-\frac76
<
-\frac87.
\]

Чтобы быстро находить дополнительные множители, разложим общий знаменатель:

\[
84=2^2\cdot3\cdot7.
\]

Для знаменателя $7$ остаются множители

\[
2^2\cdot3=12.
\]

Для знаменателя $6=2\cdot3$ остаются

\[
2\cdot7=14.
\]

Для знаменателя $4=2^2$ остаются

\[
3\cdot7=21.
\]

\warning{
Среди отрицательных чисел меньше то число, модуль которого больше:
\[
-105<-98<-96.
\]
}

% ---------------------------------------------------------

\sectiontitle{6. Вынесение общего множителя}

Если в нескольких слагаемых повторяется один и тот же множитель,
его можно вынести за скобки:

\[
ab+ac=a(b+c).
\]

Этот приём часто значительно упрощает вычисления.

Например,

\[
47{,}524\cdot65
+
47{,}524\cdot35.
\]

Выносим общий множитель:

\[
47{,}524(65+35).
\]

Получаем

\[
47{,}524\cdot100
=
4752{,}4.
\]

\keyidea{
Перед длинными вычислениями ищите общий множитель и удобную сумму
внутри скобок.
}

При умножении десятичной дроби на $10$, $100$, $1000$
запятая сдвигается вправо соответственно на $1$, $2$, $3$ цифры.

Например,

\[
47{,}524\cdot10=475{,}24,
\]

\[
47{,}524\cdot100=4752{,}4,
\]

\[
47{,}524\cdot1000=47524.
\]

% ---------------------------------------------------------

\sectiontitle{7. Простейшие неравенства и числовая прямая}

Рассмотрим неравенство

\[
x\ge -14{,}7.
\]

Оно означает: подходят все числа, которые больше либо равны $-14{,}7$.

\begin{center}
\begin{tikzpicture}[x=0.72cm,y=1cm]

  \draw[->,line width=0.9pt]
    (-3.5,0)--(8.3,0);

  \foreach \x/\lab in {
    -2/-16,
    0/-15,
    2/-14,
    4/-13,
    6/-12,
    8/-11
  }{
    \draw (\x,0.11)--(\x,-0.11);
    \node[below] at (\x,-0.11) {$\lab$};
  }

  \coordinate (P) at (0.6,0);

  \fill[accent] (P) circle (2.2pt);

  \node[above,text=accent] at (0.6,0.12)
    {$-14{,}7$};

  \draw[->,accent,line width=1.6pt]
    (0.6,0)--(8.0,0);

\end{tikzpicture}
\end{center}

Целые числа, удовлетворяющие этому неравенству:

\[
-14,-13,-12,\ldots
\]

Следовательно, наименьшее целое решение равно

\[
\boxed{-14}.
\]

\keyidea{
Для $x\ge a$ рассматривайте числа справа от $a$ на числовой прямой.
Наименьшее целое решение --- первое целое число, которое попало
в подходящую область.
}

\warning{
Знак $\ge$ читается «больше или равно»,
знак $\le$ --- «меньше или равно».

При умножении или делении обеих частей неравенства
на отрицательное число знак неравенства меняется
на противоположный.
}

% ---------------------------------------------------------

\sectiontitle{8. Короткий алгоритм перед ответом}

\begin{enumerate}

  \item
  Определите, какая информация дана:
  значения букв, значение целого выражения,
  рисунок, проценты, дроби или неравенство.

  \item
  Переведите условие задачи на математический язык.

  \item
  Выберите подходящий приём:
  подстановка, пропорция, общий знаменатель,
  вынесение множителя или числовая прямая.

  \item
  Проверьте знаки и ограничения.

  \item
  Выполните вычисления.

  \item
  Сопоставьте полученный результат с вопросом задачи.

  \item
  В тестовом задании запишите номер верного варианта,
  если именно этого требует формат работы.

\end{enumerate}

% =========================================================
% ТРЕНИРОВОЧНЫЙ БЛОК
% =========================================================

\clearpage

\sectiontitle{Тренировочный блок · Путь А}

\textbf{10 упражнений: 3 средних, 6 выше среднего, 1 сложное.}

Решения оформляйте последовательно.
Там, где это помогает рассуждению, используйте рисунок
или числовую прямую.

\begin{enumerate}

% ---------- 1 ----------

\item
\textbf{Средний уровень.}

Найдите значение выражения

\[
18a-5b+3c,
\]

если

\[
a=0{,}4,
\qquad
b=2{,}6,
\qquad
c=1{,}5.
\]

% ---------- 2 ----------

\item
\textbf{Средний уровень.}

Известно, что

\[
4mn-3k=8.
\]

Найдите значение выражения

\[
\frac{6}{4mn-3k}.
\]

% ---------- 3 ----------

\item
\textbf{Средний уровень.}

Найдите число, $12\%$ которого равны $54$.

% ---------- 4 ----------

\item
\textbf{Выше среднего.}

Квадрат имеет сторону $a$ см.
От него отрезали прямоугольную полосу
со сторонами $a$ см и $b$ см, где

\[
0<b<a.
\]

Запишите площадь оставшейся части:

\begin{enumerate}

  \item
  в виде произведения;

  \item
  после раскрытия скобок.

\end{enumerate}

% ---------- 5 ----------

\item
\textbf{Выше среднего.}

Вычислите рациональным способом:

\[
36{,}842\cdot27
+
36{,}842\cdot73.
\]

% ---------- 6 ----------

\item
\textbf{Выше среднего.}

Расположите в порядке возрастания числа

\[
-\frac{11}{12},
\qquad
-\frac56,
\qquad
-\frac78.
\]

Используйте общий знаменатель.

% ---------- 7 ----------

\item
\textbf{Выше среднего.}

Укажите наименьшее целое число,
удовлетворяющее неравенству

\[
x\ge -8{,}4.
\]

Для самопроверки изобразите ситуацию
на числовой прямой.

% ---------- 8 ----------

\item
\textbf{Выше среднего.}

После увеличения неизвестного числа на $15\%$
получили $690$.

Найдите исходное число.

% ---------- 9 ----------

\item
\textbf{Выше среднего.}

Известно, что

\[
2x-5y=-6.
\]

Найдите значение выражения

\[
7-3(2x-5y).
\]

% ---------- 10 ----------

\item
\textbf{Сложное.}

Найдите наименьшее целое $x$,
удовлетворяющее неравенству

\[
\frac{x}{3}-2{,}4\ge -5{,}1.
\]

Затем найдите значение выражения

\[
2x^2-5x+1
\]

при найденном значении $x$.

\end{enumerate}

% =========================================================
% ОТВЕТЫ
% =========================================================

\clearpage

\sectiontitle{Ответы к тренировочному блоку}

\renewcommand{\arraystretch}{1.35}

\begin{table}[ht]
\centering

\caption{Краткие ответы}

\begin{tabularx}{\linewidth}{
  @{}
  >{\bfseries}p{1.2cm}
  X
  @{}
}

\toprule

№ & Ответ \\

\midrule

1
&
$-1{,}3$
\\

2
&
$\dfrac34=0{,}75$
\\

3
&
$450$
\\

4
&
1) $a(a-b)$;
\quad
2) $a^2-ab$
\\

5
&
$3684{,}2$
\\

6
&
$
-\dfrac{11}{12}
<
-\dfrac78
<
-\dfrac56
$
\\

7
&
$-8$
\\

8
&
$600$
\\

9
&
$25$
\\

10
&
$x=-8$;
\quad
$2x^2-5x+1=169$
\\

\bottomrule

\end{tabularx}
\end{table}

\end{document}